\subsection{交换群上的矩阵}

设 G 是一个交换群（写成加法）。带给定指标集 I, K 的 G 上矩阵的集合具有交换群结构，因为它是 \(I \times K\) 到 G 的映射集合；这个群写成加法，于是若 \(M = (m_{\iota\kappa})\) 与 \(M' = (m_{\iota\kappa}')\) 是它的两个元素，则

\[
M + M ^ {\prime} = (m _ {\iota \kappa} + m _ {\iota \kappa} ^ {\prime});
\]

因此，这个群的单位元是所有元素均为零的矩阵（称为零矩阵）。显然，

\[
{ } ^ { t } ( M + M ^ { \prime } ) = { } ^ { t } M + { } ^ { t } M ^ { \prime } .\tag{2}
\]

两个矩阵的和仅在这两个矩阵的行索引集和列索引集相同时才有定义。

设 \(\mathbf{H}'\) , \(\mathbf{H}''\) 是两个集合， \(\mathbf{G}\) 是一个交换群（写成加法），而 \(f: (h', h'') \mapsto h'h''\) 是 \(\mathbf{H}' \times \mathbf{H}''\) 到 \(\mathbf{G}\) 的一个映射。给定两个矩阵

\[
M ^ {\prime} = \left(m _ {i k} ^ {\prime}\right) _ {(i, k) \in I \times K}, \quad M ^ {\prime \prime} = \left(m _ {k l} ^ {\prime \prime}\right) _ {(k, l) \in K \times L}
\]

分别位于 H' 与 H'' 上，使得 M' 的列的指标集 K 是有限的且等于 M'' 的行的指标集，则 M' 与 M'' 经由 f 的乘积，记作 M'M'' 或 f(M', M'')，是矩阵

\[
\left(\sum_ {k \in \mathbf {K}} m _ {i k} ^ {\prime} m _ {k l} ^ {\prime \prime}\right) _ {(i, l) \in \mathbf {I} \times \mathbf {L}}\tag{3}
\]

（ G 上的）。

上述定义假定 \(M'\) 的列的指标集等于 \(M''\) 的行的指标集；特别地，乘积 \(M''M'\) 没有意义，若 \(I \neq L\) 。在公式 (3) 中，同一

 \(M'\) 的同一行的元素以 \(M''\) 的同一列的元素右乘而出现；称这种乘法是按“行乘列”进行的。

设 \(f^0\) 是映射 \((h'', h') \mapsto h'h''\) （从 \(H'' \times H'\) 到 G）；由定义立即可得

\[
{ } ^ { t } ( M ^ { \prime } M ^ { \prime \prime } ) = { } ^ { t } M ^ { \prime \prime } . { } ^ { t } M ^ { \prime }\tag{4}
\]

其中左边（相应地右边）的乘积是经由 \(f\) （相应地经由 \(f^0\) ）计算的。

当 \(\mathbf{H}'\) 与 \(\mathbf{H}''\) 本身是交换群（写成加法）且 \(f\) 是 \(\mathbf{Z}\) -双线性（ \(\S 3\) ，第 1 号）时，分配律公式

\[
\left\{ \begin{array}{l} (M ^ {\prime} + N ^ {\prime}) M ^ {\prime \prime} = M ^ {\prime} M ^ {\prime \prime} + N ^ {\prime} M ^ {\prime \prime} \\ M ^ {\prime} (M ^ {\prime \prime} + N ^ {\prime \prime}) = M ^ {\prime} M ^ {\prime \prime} + M ^ {\prime} N ^ {\prime \prime} \end{array} \right.\tag{5}
\]

立即得到验证，指标集使得其中出现的和与积均有定义。

现在设 \(H_{1}\) , \(H_{2}\) , \(H_{3}\) , \(H_{12}\) , \(H_{23}\) 与 H 是交换群（写成加法）， \(f_{12}:H_{1} \times H_{2} \to H_{12}\) , \(f_{23}:H_{2} \times H_{3} \to H_{23}\) 是映射，而

\[
f _ {3}: \mathrm{H} _ {1 2} \times \mathrm{H} _ {3} \rightarrow \mathrm{H}, \quad f _ {1}: \mathrm{H} _ {1} \times \mathrm{H} _ {2 3} \rightarrow \mathrm{H}
\]

Z-双线性映射；进一步假设，对所有 \(x_{i} \in H_{i}\) （i = 1, 2, 3）

\[
f _ {3} (f _ {1 2} (x _ {1}, x _ {2}), x _ {3}) = f _ {1} (x _ {1}, f _ {2 3} (x _ {2}, x _ {3}))
\]

（它也可像上面那样写成 \((x_{1}x_{2})x_{3} = x_{1}(x_{2}x_{3}))\) ；于是，若 \(M' = (m'_{rs})\) , \(M'' = (m''_{st})\) , \(M''' = (m''_{tu})\) 分别是 \(\mathrm{H}_1\) , \(\mathrm{H}_2\) , \(\mathrm{H}_3\) 上的矩阵，

\[
(M ^ {\prime} M ^ {\prime \prime}) M ^ {\prime \prime \prime} = M ^ {\prime} (M ^ {\prime \prime} M ^ {\prime \prime \prime})\tag{6}
\]

（当两边的乘积（分别经由 \(f_{12}\) , \(f_{3}\) , \(f_{23}\) 与 \(f_{1}\) 计算）都有定义时）成立；因为

\[
\begin{array}{r l} \sum_ {t} \left(\sum_ {s} m _ {r s} ^ {\prime} m _ {s t} ^ {\prime \prime}\right) m _ {t u} ^ {\prime \prime \prime} & = \sum_ {t} \sum_ {s} (m _ {r s} ^ {\prime} m _ {s t} ^ {\prime \prime}) m _ {t u} ^ {\prime \prime \prime} = \sum_ {s} \sum_ {t} m _ {r s} ^ {\prime} (m _ {s t} ^ {\prime \prime} m _ {t u} ^ {\prime \prime \prime}) \\ & = \sum_ {s} m _ {r s} ^ {\prime} \left(\sum_ {t} m _ {s t} ^ {\prime \prime} m _ {t u} ^ {\prime \prime \prime}\right) \end{array}
\]

根据所作的假设，

（6）式的两边也记作 \(M'M''M'''\) 。对于多于三个因子的乘积，也采用类似的约定。

注记。上述公式可推广到更一般的情形。确切地说：

\begin{enumerate}
\def\labelenumi{(\alph{enumi})}
\item
假设 \(H = \bigcup_{(\iota, \kappa) \in I \times K} G_{\iota\kappa}\) ，其中每个 \(G_{\iota\kappa}\) 是写成加法的交换群；那么和 \(M + M'\) 可以在如下情形定义：对每个有序对 \((\iota, \kappa)\) ， \(m_{\iota\kappa} \in G_{\iota\kappa}\) 且 \(m'_{\iota\kappa} \in G_{\iota\kappa}\) 。
\item
  设 I、K、L 为三个集合，其中 K 假定为有限集，且设 \(H' = \bigcup_{(i, k) \in I \times K} H'_{ik}\) , \(H'' = \bigcup_{(k, l) \in K \times L} H''_{kl}\) , \(H = \bigcup_{(i, l) \in I \times L} H_{il}\) 是三个集合；假设每一个 \(H_{il}\) 是一个交换群，以加法形式书写，并且对于每个三元组 \((i, k, l)\) ，设
\end{enumerate}

\[
f _ {i k l}: \mathrm{H} _ {i k} ^ {\prime} \times \mathrm{H} _ {k l} ^ {\prime \prime} \rightarrow \mathrm{H} _ {i l}
\]

是一个映射。那么，若 \(M' = (m_{ik}')_{(i,k) \in I \times K}\) , \(M'' = (m_{kl}'')_{(k,l) \in K \times L}\) 是满足 \(m_{ik}' \in H_{ik}'\) 与 \(m_{kl}'' \in H_{kl}''\) （对所有 \(i, k, l\) ）的矩阵，我们就可以定义乘积 \(M'M''\) （经由诸 \(f_{ikl}\) ）。我们把写出并证明类似于 (4)、(5) 与 (6) 的公式的工作留给读者。

